% !TEX root = ../main.tex
% =====================================================================
% §1 Introduction — gdn2vessel benchmark-only D&B 稿
% 真源：STORY_FRAMEWORK.md（2026-06-24 benchmark-only 重定位）+ LEADERBOARD_MATRIX.md
% 承重点 = benchmark 能区分方法（hedge 至批2，禁写死 sharply ranks）+ 诊断现有方法/指标局限。
% 数字一律 \TODO 占位（禁臆造）；引用 \cite 占位（refs.bib 待建）。双盲：无自引/机构。
% =====================================================================
\section{Introduction}
\label{sec:intro}

\subsection{Reconnection at vessel breakpoints is clinically decisive yet hard to score}
\label{sec:intro:problem}

Thin tubular structures---retinal vessels, airways, and neuronal
processes---carry information whose clinical value lies primarily in their
\emph{connectivity}: whether a vessel tree remains a single connected
component, how branches join, and where a path is interrupted. In 2D retinal
fundus imaging, automated segmentation routinely produces masks that are
locally accurate yet topologically broken: a vessel that is faint under low
contrast, crossed by another vessel, or partially occluded is frequently
rendered as two disconnected fragments rather than one continuous path.
These \emph{breakpoints} matter downstream, because vessel tortuosity,
branching geometry, and the integrity of the arteriovenous network are read
off the connected topology rather than off per-pixel overlap
\cite{TODO_clinical_connectivity}.

A model can therefore score well on conventional overlap metrics
(Dice, IoU) while failing precisely where it matters: at the gaps. The
ability to \emph{reconnect} a vessel across such a gap is a distinct skill
from segmenting it, and the central difficulty addressed in this paper is
not how to perform reconnection, but how to \emph{measure} it. Reconnection
quality is hard to score for two reasons. First, overlap metrics are blind
to connectivity by construction, so a method that closes a gap and a method
that leaves it open can receive almost identical Dice. Second, the
connectivity-aware metrics that do exist---and the bespoke ``success''
measures proposed alongside individual reconnection methods---are reported
under heterogeneous, often unspecified protocols, which makes cross-method
comparison unreliable and leaves their own failure modes undocumented.

\subsection{There is no standardized 2D-fundus reconnection evaluation suite}
\label{sec:intro:gap}

Progress on a perceptual capability is usually paced by a standardized
benchmark with a fixed protocol, a leaderboard, and metrics whose semantics
are understood. Vessel reconnection has the methods but not the benchmark.
On the method side, dedicated reconnection and connectivity-recovery
approaches exist---for example region-aware co-segmentation for vessel
reconnection (CorSegRec) \cite{TODO_corsegrec} and global--local
connectivity priors (GLCP) \cite{TODO_glcp}---and 2D-fundus completion has
been studied as a generative task (MaskVSC) \cite{TODO_maskvsc}. These are
\emph{methods}, not evaluation suites, and each reports connectivity under
its own protocol. On the benchmark side, a synthetic-breakpoint benchmark
already exists, but in a different setting: PTR \cite{TODO_ptr} provides a
standardized disconnection benchmark for \emph{3D pulmonary} airways/vessels,
not 2D fundus. We therefore make \emph{no} claim to the first disconnection
benchmark; PTR occupies that ground in 3D lung.

The gap we identify is specific. For \emph{2D retinal fundus} reconnection
there is, to our knowledge, no standardized suite that (i) injects
breakpoints under a fixed, reproducible protocol, (ii) reports a coherent
family of connectivity metrics whose failure modes are characterized, and
(iii) places a broad set of segmentation and reconnection methods on a
single comparable leaderboard. Without such a suite, reported reconnection
numbers are not comparable across papers, and the metrics themselves are
used without knowing whether they can be exploited. The latter point is not
hypothetical: as we show, one widely intuitive ``did the gap close?''
success measure can be saturated by a degenerate over-predicting model.

\subsection{A 2D-fundus disconnection--reconnection evaluation suite}
\label{sec:intro:hook}

We present, to our knowledge, the first standardized evaluation suite for
\textbf{2D retinal-fundus vessel reconnection}. The suite combines four
components that, taken together, constitute its novelty---no single piece is
claimed as first in isolation:

\begin{itemize}
  \item a \textbf{synthetic-breakpoint protocol} that injects gaps with
        fixed radius and gap parameters, aligned with a published
        plug-and-play breakpoint procedure \cite{TODO_creatis} (whose method
        paper we use as a protocol reference, not a competitor), arranged
        into a four-level \emph{severity} grid (Easy / Medium / Hard /
        Extreme); evaluation sets are strictly held out, and the protocol
        never reads ground-truth topology into any evaluated method;
  \item a \textbf{family of connectivity metrics} for reconnection,
        cross-checked against standard topological measures (clDice, Betti
        errors), \emph{together with a documented diagnosis of a metric
        failure mode}: a ``gap closed'' success rate that we show can be
        saturated by over-prediction;
  \item a \textbf{full-spectrum leaderboard} that scores a broad set of
        segmentation and reconnection methods on the same protocol across
        datasets and severities, with \emph{severity-response} profiles that
        track how reconnection quality degrades as breakpoints become harder;
  \item a \textbf{pre-registered diagnostic methodology} with kill-criteria
        fixed before the data were seen, and the \emph{honest negative
        findings} it produced.
\end{itemize}

The load-bearing contribution is that the suite \emph{can distinguish methods
and can expose limitations of existing methods and metrics}---not that any
particular method wins. We are explicit about the status of these claims.
The negative findings and the metric-saturation diagnosis are established
results that we report as such. The strength of the separation the
benchmark induces between methods, however, is reported as what the
benchmark reveals rather than asserted as a sharp ranking: we present the
leaderboard and the severity-response profiles and report the separation the
benchmark exhibits, without claiming that it sharply ranks methods
\TODO{批2 区分度命门未跑——确认后据实更新 hedge 措辞}.

\subsection{Contributions}
\label{sec:intro:contrib}

\begin{enumerate}
  \item \textbf{A benchmark and synthetic-breakpoint protocol for 2D
        retinal-fundus reconnection.} We build a reproducible, held-out
        evaluation suite with a fixed-parameter breakpoint-injection
        protocol arranged over a four-level severity grid. We deliberately
        scope the novelty to ``the first \emph{2D-fundus} reconnection
        evaluation suite as a combination'' rather than the first
        disconnection benchmark, which PTR \cite{TODO_ptr} already provides
        in 3D lung.

  \item \textbf{A connectivity-metric family with a documented
        over-prediction failure mode.} We report a family of reconnection
        metrics cross-checked against clDice and Betti errors, and we
        diagnose that the ``gap closed'' success rate can be saturated by
        over-prediction: a degenerate model with Dice~$\approx$~\TODO{0.16}
        attains a success rate of \TODO{1.0}, because the underlying gap-closure
        test only checks whether any foreground pixel covers the centre of
        the erased region. We therefore use a topological continuity measure
        as the primary reconnection axis and report the success rate only
        alongside a specificity / over-prediction context. This metric
        diagnosis follows the tradition of metric-failure-mode analysis
        \cite{TODO_pixelwise,TODO_touchstone}.

  \item \textbf{A full-spectrum leaderboard with severity-response curves.}
        We evaluate \TODO{$\geq$12} segmentation and reconnection baselines on a
        single protocol across datasets and four severity levels, reporting
        a segmentation sanity axis, a topology axis, and a reconnection axis,
        and we characterize how reconnection quality degrades with breakpoint
        severity. We report the separation the benchmark reveals and the
        severity-response profiles, deferring any stronger ranking claim
        \TODO{批2 区分度命门未跑——核 verifier 后更新}.

  \item \textbf{A pre-registered diagnostic methodology and honest negative
        findings.} We treat two of the project's load-bearing hypotheses as
        examples of falsify-the-crux discipline, with kill-criteria fixed
        before evaluation. The pre-registered audits show that a delta-rule
        associative-memory module is \emph{not} mechanistically specific on
        2D vessels (it matches a plain gated-memory baseline at the only
        comparable converged operating point \TODO{MQAR n=64, lr 1e-3:
        gdn2 vs.\ gla, 核 csv}), and that a differentiable vesselness gate
        does \emph{not} robustly improve topology (its pre-registered gate
        failed: positive on one dataset, negative on another). We frame these
        negative findings as positive evidence that the benchmark has
        discriminative power: a suite that can show a gated variant to be
        unstable and a memory mechanism to be non-specific is doing its job.
        The corresponding modules appear on the leaderboard as evaluated
        methods, not as headline contributions.
\end{enumerate}

\noindent The remainder of the paper details the benchmark protocol and
connectivity-metric family (Section~\ref{sec:benchmark}), the baselines and
leaderboard, and the diagnostic methodology and negative findings.
